\documentclass[a4paper,12pt]{article}
\setlength{\headheight}{68.803pt}
\addtolength{\topmargin}{-56.803pt}

\usepackage[utf8]{inputenc}
\usepackage[T1]{fontenc}
\usepackage{lmodern}
\usepackage[nopatch=footnote]{microtype}
\usepackage{fvextra}
\DefineVerbatimEnvironment{Verbatim}{Verbatim}{
  breaklines=true,
  breakanywhere=true,
  frame=single,
  fontsize=\small
}

% margin settings
\usepackage[
  left=1in,
  right=1in,
  top=3cm,
  bottom=1in
]{geometry}

\usepackage[
    colorlinks=true,
    linkcolor=black,
    urlcolor=blue,
    citecolor=blue,
    pdfborder={0 0 0},
    pdfauthor={Massimo Mantineo},
    pdftitle={Server TCP Bloccante e Non-Bloccante},
    pdfsubject={Laboratorio di Reti e Sistemi Distribuiti: HANDS-ON 2},
]{hyperref}

\usepackage{graphicx}
\graphicspath{{../../assets/}}
\usepackage{tikz}
\usetikzlibrary{positioning, arrows.meta}
\usepackage{listings}
\usepackage{xcolor}
\usepackage{amsmath, amssymb}
\usepackage{fancyhdr}
\usepackage{setspace}
\usepackage{pgfplots}
\pgfplotsset{compat=1.17}
\usepackage{booktabs}
\usepackage[skins, listings]{tcolorbox}
\usepackage{float}

\newtcblisting{terminal}[1][]{
    listing only,
    colback=black!85!white,
    colframe=black!70!white,
    colupper=green!80!white,
    coltitle=white,
    title=#1,
    fonttitle=\bfseries\sffamily,
    listing options={
        language={},
        basicstyle=\ttfamily\small\color{green!80!white},
        backgroundcolor={},
        frame=none,
        breaklines=true,
        showstringspaces=false,
        numbers=none
    },
    enhanced,
    drop shadow,
    arc=3mm
}

\newtcblisting{ccode}[1][]{
    listing only,
    colback=blue!3!white,
    colframe=blue!70!black,
    title=#1,
    fonttitle=\bfseries\sffamily,
    listing options={
        style=cstyle,
        backgroundcolor={},
        frame=none
    },
    enhanced,
    drop shadow,
    arc=3mm
}

\setlength{\footskip}{50pt}

% settings for C listings
\lstdefinestyle{cstyle}{
    language=C,
    basicstyle=\ttfamily\small,
    numbers=left,
    numberstyle=\tiny,
    stepnumber=1,
    numbersep=5pt,
    backgroundcolor=\color{white},
    showspaces=false,
    showstringspaces=false,
    showtabs=false,
    frame=single,
    tabsize=4,
    captionpos=b,
    breaklines=true,
    breakatwhitespace=true,
    keywordstyle=\color{blue}\bfseries,
    commentstyle=\color{green!50!black},
    stringstyle=\color{red},
}
\lstset{style=cstyle}

% Header and Footer
\pagestyle{fancy}
\fancyhf{}

\fancyhead[L]{%
  \parbox[b]{0.45\textwidth}{%
    \small
    Student Name: Massimo Mantineo\\
    Student ID: 541924
    \vspace{20pt}
  }%
}
\fancyhead[C]{%
  \parbox[b]{0.2\textwidth}{%
    \centering
    \normalsize \bfseries
    LabReSiD25\\
    Hands-On 2
    \vspace{20pt}
  }%
}
\fancyhead[R]{%
  \parbox[b]{0.35\textwidth}{%
    \flushright
    \includegraphics[width=3.5cm]{\detokenize{logo_unime_orizontale.png}}
    \vspace{15pt}
  }%
}

\renewcommand{\contentsname}{Indice}

\title{\vspace{-2cm}Relazione: Server TCP Bloccante e Non-Bloccante\\
\large (Hands-On 2)}
\author{Massimo Mantineo \\ Matricola: 541924}
\date{MESSINA, A.A. 2024/2025}

\begin{document}

% Cover page
\begin{titlepage}
    \thispagestyle{empty}
    \centering
    \vspace*{1cm}
    \includegraphics[width=6cm]{\detokenize{logo_unime_orizontale.png}}\par\vspace{1cm}
    \vspace{2cm}
    {\Large \textbf{Università degli Studi di Messina}\par}
    {\large Dipartimento MIFT - Corso di Laurea Triennale in Informatica (L-31)\par}
    \vspace*{1cm}
    {\large Laboratorio di Reti e Sistemi Distribuiti\par}
    \vspace{1cm}
    {\Large \textbf{HANDS-ON 2:}\\
    Server TCP Bloccante/Non-Bloccante\par}
    \vspace{2cm}
    {\large Massimo Mantineo \par}
    {\large Matricola: 541924 \par}
    \vfill
    {\large Messina, A.A. 2024/2025 \par}
\end{titlepage}

\clearpage
\pagestyle{fancy}
\tableofcontents
\newpage

\section{Problem Statement}
L'obiettivo della seconda esercitazione (Hands-On 2) è l'analisi comparativa dei modelli di I/O nel contesto della programmazione di rete sotto Linux. Il \textit{core} dell'attività consiste nell'implementazione di un servizio di validazione della \textbf{Successione di Fibonacci}. La sfida tecnica risiede nel gestire flussi di dati sequenziali in due scenari distinti: un server a modello \textbf{bloccante} (iterativo) e un server \textbf{non-bloccante} basato su multiplexing. Il sistema deve garantire che ogni numero ricevuto dal client segua rigorosamente la legge di ricorrenza $F_n = F_{n-1} + F_{n-2}$, restituendo un acknowledgment (\texttt{ACK_FIB_OK}) o segnali di errore in caso di divergenza.

\section{Stato dell'Arte}
La gestione della concorrenza e delle attese è un pilastro dei sistemi distribuiti.
\begin{itemize}
    \item \textbf{I/O Bloccante}: Rappresenta il comportamento di default dei socket Unix. Una chiamata a \texttt{read()} sospende il thread del chiamante finché il buffer di ricezione del kernel non contiene dati. Sebbene semplice, questo modello soffre di gravi problemi di scalabilità (scarsa \textit{responsiveness}) in presenza di client lenti o multipli.
    \item \textbf{I/O Non-Bloccante}: Tramite la syscall \texttt{fcntl(O_NONBLOCK)}, il socket viene istruito a non sospendere il processo; se i dati non sono pronti, la funzione ritorna immediatamente l'errore \texttt{EAGAIN}.
    \item \textbf{Multiplexing con \texttt{select()}}: Per evitare il \textit{busy-waiting} (polling continuo) su socket non-bloccanti, si utilizza la \texttt{select()}, che permette di monitorare un intero set di file descriptor, svegliando il processo solo quando almeno uno di essi è pronto per un'operazione di I/O.
\end{itemize}

\section{Metodologia ed Implementazione}
L'implementazione prevede un client generatore e due varianti di server. La logica di validazione della sequenza di Fibonacci è il cuore algoritmico del sistema.

\subsection{Client: Generazione Dinamica}
Il client calcola i termini della successione e li trasmette serializzati come stringhe.
\begin{ccode}[Logica Fibonacci Client]
long long a = 0, b = 1, next;
for (int i = 0; i < MAX_FIB; i++) {
    sprintf(msg, "%lld", a);
    send(sock, msg, strlen(msg), 0);
    // ... attesa ACK ...
    next = a + b; a = b; b = next;
}
\end{ccode}

\subsection{Server: Modelli a Confronto}
Il server bloccante utilizza un semplice loop \texttt{accept} $\rightarrow$ \texttt{read}. Il server non-bloccante, invece, configura i socket tramite:
\begin{ccode}[Snippet Configurazione Non-Bloccante]
int flags = fcntl(master_socket, F_GETFL, 0);
fcntl(master_socket, F_SETFL, flags | O_NONBLOCK);
activity = select(max_sd + 1, &readfds, NULL, NULL, NULL);
\end{ccode}
Questa architettura permette di mantenere stati di validazione indipendenti per ogni client tramite una struttura \texttt{client_state}.

\section{Risultati}
La validazione è stata condotta simulando accessi concorrenti e ispezionando l'interlacciamento dei pacchetti.

\subsection{Validazione Funzionale}
Entrambi i server riconoscono correttamente la sequenza di Fibonacci. Tuttavia, i log evidenziano la differenza nel determinismo temporale:
\begin{terminal}[Log Multiplexing]
New connection, socket fd: 4, ip: 127.0.0.1
Received: 0 from socket 4 -> ACK_FIB_OK
Received: 1 from socket 5 -> ACK_FIB_OK (Gestione Parallela)
\end{terminal}

\subsection{Analisi Wireshark}
L'analisi del traffico conferma che nel modello bloccante un secondo client rimane in attesa (nonostante l'handshake kernel sia avvenuto) finché il primo non chiude la sessione. Nel modello non-bloccante, i pacchetti \texttt{PSH, ACK} di diversi client risultano mescolati cronologicamente, dimostrando il successo del multiplexing.

\section{Conclusioni e Sviluppi Futuri}
L'Hands-On 2 ha dimostrato che il multiplexing dell'I/O è essenziale per la creazione di server reattivi. La gestione manuale degli stati dei client (\textit{state machines}) è un compromesso necessario per ottenere prestazioni elevate in un singolo thread.

\noindent \textbf{Sviluppi futuri:}
\begin{itemize}
    \item \textbf{Passaggio a Epoll}: Evolvere verso \texttt{epoll()} per gestire migliaia di connessioni simultanee con complessità $O(1)$ invece di $O(N)$ della \texttt{select()}.
    \item \textbf{Load Balancing}: Implementazione di una coda di messaggi per distribuire il carico di validazione computazionale su un pool di worker thread.
\end{itemize}

\end{document}
